% =====================================================================
\section{Discussion}\label{sec:discussion}
% =====================================================================
\subsection{Main findings}
On a held-out IDRiD test set that played no role in any modelling decision, NITSCI reached a five-class accuracy of
\res{acc5}\% with a QWK of \res{qwk5}, graded \res{within1}\% of images within one level of the reference, and
detected referable DR with an AUC of \res{refauc} at a sensitivity of \res{refsens}\% and a specificity of
\res{refspec}\%. \authnote{Summarise in two or three sentences, using only claims supported by
Tables~\ref{tab:main}--\ref{tab:infer}: (1) whether the full model improves on the plain baseline beyond the bootstrap
interval; (2) which ablations had the largest effect; (3) whether the ordinal head and threshold fitting reduced large
errors, visible in QWK and within-one-grade accuracy.}

\subsection{Interpretation of the components}
\paragraph{Cross-dataset pretraining.} APTOS-2019 differs from IDRiD in cameras, image quality and population, but
shares the grading scale and the visual vocabulary of DR lesions. Pretraining on it exposes the model to roughly eight
times more graded fundus images before it sees IDRiD, which is consistent with the general observation that
in-domain data narrows the gap left by ImageNet pretraining \citep{raghu2019transfusion}. \authnote{Report the size of
the ``w/o APTOS'' effect from Table~\ref{tab:ablation} and whether it is the largest single factor.}

\paragraph{Lesion expert, cross-attention and gate.} The three architectural components address the same problem from
different angles: the lesion map concentrates the representation on a few informative locations, cross-attention
lets lesion tokens be interpreted in their retinal context, and the gate allows the balance between global and lesion
evidence to vary from image to image. A plausible expectation is that the gate favours lesion evidence for borderline
images between grades 0, 1 and 2, where the decision depends on a few lesions, and global context for advanced
disease, where lesions are widespread. \authnote{Check this with the gate values exported for the test images, or
remove the sentence. Report which of the three ablations changed accuracy or QWK beyond the confidence interval,
and say plainly if a component had no measurable effect at this test-set size.}

\paragraph{Ordinal head and decision rule.} Combining the expected grade of the classifier with a regression output
produces a score that changes smoothly with disease severity, so errors tend to fall on adjacent grades. Fitting the
cut-points on out-of-fold predictions adapts them to the class balance of IDRiD without touching the test set. The
comparison of decision rules in Table~\ref{tab:infer} shows how much of the final accuracy is due to the rule rather
than to the network. \authnote{Comment on whether the OOF ranking of the rules matched their test ranking; if not,
this indicates that the threshold fit partly reflects noise in a small OOF set.}

\paragraph{Ensembling and TTA.} Averaging 20 predictions per image reduces variance, which matters most for the
under-represented grades. The difference between a single fold model and the ensemble (Table~\ref{tab:infer}) also
quantifies a practical trade-off: a single model without TTA is twenty times cheaper at inference.

\subsection{Comparison with published results}
Direct numerical comparison with the IDRiD grading challenge \citep{porwal2020idrid} is not possible from this study,
because the challenge uses an official train/test partition whereas we use a stratified random hold-out of the Kaggle
copy. Numbers reported on different splits, with different pre-processing and sometimes with thresholds tuned on the
test set, are not comparable, and placing them in the same table would mislead the reader. We therefore compare
against a plain baseline and ablations trained under identical conditions (Table~\ref{tab:ablation}).
\authnote{Either (a) re-run the pipeline on the official IDRiD grading partition and add a table comparing with the
challenge results reported in \citet{porwal2020idrid}, copying the numbers from that paper; or (b) keep this paragraph.}

\subsection{Interpretability}
The lesion-attention map and Grad-CAM \citep{selvaraju2017gradcam} offer a view of where the model looks, which can help
clinicians to spot obvious failure modes. They should not be over-interpreted: saliency maps can be insensitive to the
model parameters \citep{adebayo2018sanity}, and post-hoc explanations of black-box models are not a substitute for
validated, inherently interpretable evidence in high-stakes decisions \citep{rudin2019stop}. Because our lesion map is
learned without lesion masks, we treat it as a qualitative aid only; validating it against the IDRiD lesion
annotations is part of future work.

\subsection{Limitations}
\begin{itemize}
\item \textbf{Small test set.} With \res{ntest} test images, confidence intervals are wide, and per-grade metrics for
  rare grades rest on few images. Differences of a few percentage points between configurations may not be
  statistically detectable.
\item \textbf{Single target dataset.} All results come from IDRiD, collected at one centre with one camera.
  Performance on other populations, cameras and image-quality levels, for example Messidor
  \citep{decenciere2014messidor} or DDR \citep{li2019ddr}, has not been tested.
\item \textbf{Non-official split.} The random hold-out prevents direct comparison with the IDRiD challenge.
\item \textbf{Weakly supervised lesion map.} The lesion-attention map has not been validated against lesion masks.
\item \textbf{Training randomness.} Variance is reflected through five folds but not through repeated runs with
  different seeds.
\item \textbf{Reference standard.} Grades are taken from the dataset providers; grader variability, which is
  substantial for DR \citep{krause2018grader}, is not modelled.
\item \textbf{Optimistic OOF estimates.} Checkpoints are selected on the same folds used for OOF predictions, so the
  cross-validation numbers (but not the test numbers) are slightly optimistic \citep{varma2006bias}.
\end{itemize}

\subsection{Ethical and clinical considerations}
NITSCI is a research prototype, not a medical device. Both datasets are public and de-identified. Several issues must
be addressed before any clinical use. First, models trained on data from one population can perform worse, or fail
silently, on others; algorithmic bias in health care has had measurable consequences \citep{obermeyer2019dissecting},
so subgroup performance by age, sex, ethnicity and camera should be reported. Second, the operating point should be
set together with clinicians to reach a target sensitivity rather than to maximise accuracy, since a false negative
delays treatment while a false positive costs a specialist visit. Third, probabilities from deep networks are often
mis-calibrated \citep{guo2017calibration}; calibration should be checked before a probability threshold is used for
referral. Finally, prospective external validation reported according to TRIPOD+AI \citep{collins2024tripod} is
required, as autonomous DR systems in clinical use have shown \citep{abramoff2018pivotal}.
\sectionend{discussion}
